\scanpage{39}

\diagram{\begin{tikzcd}[ampersand replacement=\&] X\arrow[d,"f"]\\Y=\unsure{\mathrm{Spec}|\mathrm{Spm}}(O)\end{tikzcd}}

\[
F=\widetilde M
\]

$f$ \inserted{plat et} propre, $Y$ schéma noethérien local, $F$ cohérent sur $Y$

\diagram{\begin{tikzcd}[ampersand replacement=\&,column sep=small] F'=f^*(F),\quad F'_1\subset F'\arrow[r,bend right=14] \& \text{ss faisceau cohérent.}\end{tikzcd}}

\deleted{\unsure{on}} Je dis qu'il y a un plus petit ss-faisceau cohérent $F_1$ de $F'$ t.q $f^*(F_1)\supset F'_1$

% slanted
\margin{N.B. $F$ est défini par un $O$-module $M$, les $F_i$ par des $M_i\subset M$, par \uline{définition} on pose $\unsure{\bigwedge|\bigcap}(F_i)=\widetilde{\bigcap M_i}$ [ce qui est \unsure{pas} le faisceau $\bigcap$, mais l'inf sur les faisceaux qu.coh. sur $Y$]}

Si $f_{\deleted{\unsure{*}}}^*(F_1)$ et $f^*(F_2)\supset F'_1$ alors $f_{\deleted{\unsure{*}}}^*(F_1\cap F_2)\underset{\substack{\text{platitude}\\\text{de }f}}{=}f_{\deleted{\unsure{*}}}^*(F_1)\cap f^*(F_2)\supset F'_1$, les $F_i\subset F$ tq $f^*(F_i)\supset F'_1$ \deleted{\unsure{forment}} un \deleted{filtrant} décroissant. \deleted{Il} reste à prouver que $F_1=\bigcap F_i$, alors $f^*(F_1)\supset F'_1$ i.e. l'image de $F'_1$ dans $F'/f^*(F_1)=f^*(F/F_1)$ est nulle. Remplaçant $F$ par $F/F_1$, on est ramené à prouver ceci : \uline{si} $\bigcap F_i=0$, \uline{alors} tout faisceau \deleted{\unsure{cohérent}} $F'_1$ contenu dans $\bigcap f^*(F_i)$ est nul

\margin{[est-il vrai que $\bigcap f^*(F_i)$ est nul ?]}

% slanted
\margin{Énoncé : ss les conditions envisagées, ($f$ propre, les $F_i$ décroissants dans $F$ \unsure{coh.}, $Y$ local et \unsure{noeth.}) on $\ f^*(\bigcap F_i)=\bigcap f^*(F_i)$ ?}

\deleted{$\unsure{\bigcap f^*(F_i)=0}$} ($f$ propre \deleted{\unsure{surjectif}}, les $F_i$ filtrants décr. dans $F$ cohérent, $Y$ local et noeth.). \deleted{\unsure{Comme en fait|Comme $f$} $F'_1=\unsure{\bigcap f^*}$} Comme le support de \deleted{effet $s$ une section de} $F'_1$ est fermé, il suffit de prouver qu'il ne rencontre pas $f^{-1}(y)$ ($y$ le pt fermé de $Y$) donc de prouver que si $s$ est une section de $F'_1$ sur un ouvert $U$, elle est nulle en tt $x\in f^{-1}(y)$. Or comme \deleted{\unsure{les}} $\bigcap M_i=0$, les $M_i$ tendent vers $0$ dans $M$ (?)

% slanted
\margin{\uline{Arranger les notations} !}

% archive sheet number: 38
